\documentclass[10pt,a4paper]{article}
\usepackage{amsmath,amssymb,amsthm,graphicx}
% Card-local minimal preamble for the cedram unit ahl-163-stochastic-nonlinear-wave-dynamics-on-compact-2023.
%% The journal class cedram is not on the host; the card declares the author
%% environment names itself and copies the author macro block verbatim
%% (CONTROL s8.8.166: cut the preamble, do not port class-specific lines).
\usepackage{amsmath,amssymb,amsthm,mathtools,graphicx}
\usepackage{xcolor}
\usepackage{xfrac}
\usepackage{bbm}
\newtheorem{theo}{Theorem}[section]
\newtheorem{lemm}{Lemma}[section]
\newtheorem{prop}{Proposition}[section]
\newtheorem{coro}{Corollary}[section]
\newtheorem{defi}{Definition}[section]
\newtheorem{rema}{Remark}[section]
\newtheorem{enonce}{Statement}[section]




\renewcommand{\th}{\tau_h}

\newcommand{\scrit}{s_{\textup{crit}}}
\newcommand{\sscal}{s_{\text{scal}}}
\newcommand{\sconf}{s_{\text{conf}}}
\newcommand{\damp}{\textup{damp}}




\newcommand{\I}{\hspace{0.5mm}\text{I}\hspace{0.5mm}}
\newcommand{\II}{\text{I \hspace{-2.8mm} I} }
\newcommand{\III}{\text{I \hspace{-2.9mm} I \hspace{-2.9mm} I}}
\newcommand{\IV}{\text{I \hspace{-2.9mm} V}}



\newcommand{\noi}{\noindent}

\newcommand{\Z}{\mathbb{Z}}
\newcommand{\R}{\mathbb{R}}
\newcommand{\C}{\mathbb{C}}
\newcommand{\T}{\mathbb{T}}
\newcommand{\E}{\mathbb{E}}
\newcommand{\Prob}{\mathbb{P}}
\newcommand{\N}{\mathbb{N}}


\DeclareMathOperator*{\supp}{supp}
\let\Re= \undefined
\DeclareMathOperator*{\Re}{Re}
\let\Im= \undefined
\DeclareMathOperator*{\Im}{Im}


\newcommand{\deff}{\overset{\textup{def}}{=}}


\let\P= \undefined
\newcommand{\P}{\mathbf{P}}
\newcommand{\x}{\mathbf{x}}
\newcommand{\y}{\mathbf{y}}
\newcommand{\ind}{\mathbf{1}}


\newcommand{\En}{\mathcal{E}}
\renewcommand{\L}{\mathcal{L}}
\newcommand{\A}{\mathcal{A}}
\newcommand{\F}{\mathcal{F}}
\newcommand{\B}{\mathcal{B}}
\renewcommand{\S}{\mathcal{S}}
\newcommand{\Dp}{\mathcal{D}'}
\newcommand{\J}{\mathcal{J}}
\newcommand{\M}{\mathcal{M}}
\renewcommand{\H}{\mathcal{H}}


\newcommand{\Id}{\mathrm{Id}}
\newcommand{\ZZ}{Z}
\newcommand{\crit}{\mathrm{crit}}
\newcommand{\Span}{\mathrm{Span}}
 
\newcommand{\g}{\mathfrak{g}}
\newcommand{\frkz}{\mathfrak{z}}


\newcommand{\al}{\alpha}
\newcommand{\be}{\beta}
\newcommand{\dl}{\delta}

\newcommand{\Ld}{\Lambda}
\newcommand{\ld}{\lambda}
\renewcommand{\b}{\beta}

\newcommand{\Dlg}{\Delta_{\g}}
\newcommand{\Dl}{\Delta}
\newcommand{\eps}{\varepsilon}
\newcommand{\kk}{\kappa}
\newcommand{\s}{\sigma}
\newcommand{\Si}{\Sigma}
\newcommand{\wt}{\widetilde}
\newcommand{\dx}{\partial_x}
\newcommand{\dt}{\partial_t}
\newcommand{\W}{\mathcal{W}}

\renewcommand{\o}{\omega}
\renewcommand{\O}{\Omega}



\newcommand{\les}{\lesssim}

\newcommand{\jb}[1] {\langle #1 \rangle}






\makeatletter
\def\editors#1{%
\def\editor@name{#1}
\if@francais
\def\editor@string{Recommand\'e par les \'editeurs \editor@name.}
\else
\def\editor@string{Recommended by Editors \editor@name.}
\fi}
\makeatother



%%%%%%%%%%%%%%%%%%%%%%%%%%%%%%%%%%%%%%%%%%%%%%%%%%%%%%%%


\graphicspath{{./figures/}}

\newcommand*{\mk}{\mkern -1mu}
\newcommand*{\Mk}{\mkern -2mu}
\newcommand*{\mK}{\mkern 1mu}
\newcommand*{\MK}{\mkern 2mu}

%\hypersetup{urlcolor=purple, linkcolor=blue, citecolor=red}

\newcommand*{\romanenumi}{\renewcommand*{\theenumi}{\roman{enumi}}}
\newcommand*{\Romanenumi}{\renewcommand*{\theenumi}{\Roman{enumi}}}
\newcommand*{\alphenumi}{\renewcommand*{\theenumi}{\alph{enumi}}}
\newcommand*{\Alphenumi}{\renewcommand*{\theenumi}{\Alph{enumi}}}

\usepackage{cleveref}
\crefname{originalequation}{original equation}{original equations}
\Crefname{originalequation}{Original equation}{Original equations}
\makeatletter
\newlabel{FNLW}{{5.1}{source}{}{sourcepointer-1}{}}
\newlabel{FNLW@cref}{{[originalequation][0][]5.1}{[1][0][]source}{}{}{}}
\makeatother
\crefname{originalequation}{original equation}{original equations}
\Crefname{originalequation}{Original equation}{Original equations}
\makeatletter
\newlabel{FNLW1}{{1.16}{source}{}{sourcepointer-2}{}}
\newlabel{FNLW1@cref}{{[originalequation][0][]1.16}{[1][0][]source}{}{}{}}
\makeatother
\crefname{originalequation}{original equation}{original equations}
\Crefname{originalequation}{Original equation}{Original equations}
\makeatletter
\newlabel{Hsum}{{3.1}{source}{}{sourcepointer-3}{}}
\newlabel{Hsum@cref}{{[originalequation][0][]3.1}{[1][0][]source}{}{}{}}
\makeatother
\crefname{originalequation}{original equation}{original equations}
\Crefname{originalequation}{Original equation}{Original equations}
\makeatletter
\newlabel{bd-wN}{{5.11}{source}{}{sourcepointer-4}{}}
\newlabel{bd-wN@cref}{{[originalequation][0][]5.11}{[1][0][]source}{}{}{}}
\makeatother
\crefname{originalequation}{original equation}{original equations}
\Crefname{originalequation}{Original equation}{Original equations}
\makeatletter
\newlabel{condition2}{{5.10}{source}{}{sourcepointer-5}{}}
\newlabel{condition2@cref}{{[originalequation][0][]5.10}{[1][0][]source}{}{}{}}
\makeatother
\crefname{originalequation}{original equation}{original equations}
\Crefname{originalequation}{Original equation}{Original equations}
\makeatletter
\newlabel{condition3}{{5.10}{source}{}{sourcepointer-6}{}}
\newlabel{condition3@cref}{{[originalequation][0][]5.10}{[1][0][]source}{}{}{}}
\makeatother
\crefname{originalequation}{original equation}{original equations}
\Crefname{originalequation}{Original equation}{Original equations}
\makeatletter
\newlabel{def-Psidamp}{{1.11}{source}{}{sourcepointer-7}{}}
\newlabel{def-Psidamp@cref}{{[originalequation][0][]1.11}{[1][0][]source}{}{}{}}
\makeatother
\crefname{originalequation}{original equation}{original equations}
\Crefname{originalequation}{Original equation}{Original equations}
\makeatletter
\newlabel{dlmj}{{5.9}{source}{}{sourcepointer-8}{}}
\newlabel{dlmj@cref}{{[originalequation][0][]5.9}{[1][0][]source}{}{}{}}
\makeatother
\crefname{originalequation}{original equation}{original equations}
\Crefname{originalequation}{Original equation}{Original equations}
\makeatletter
\newlabel{global-bd1}{{5.14}{source}{}{sourcepointer-9}{}}
\newlabel{global-bd1@cref}{{[originalequation][0][]5.14}{[1][0][]source}{}{}{}}
\makeatother
\crefname{originallemma}{original lemma}{original lemmas}
\Crefname{originallemma}{Original lemma}{Original lemmas}
\makeatletter
\newlabel{lem-Gibbs}{{3.2}{source}{}{sourcepointer-10}{}}
\newlabel{lem-Gibbs@cref}{{[originallemma][0][]3.2}{[1][0][]source}{}{}{}}
\makeatother
\crefname{originalproposition}{original proposition}{original propositions}
\Crefname{originalproposition}{Original proposition}{Original propositions}
\makeatletter
\newlabel{prop-FNLW}{{5.1}{source}{}{sourcepointer-11}{}}
\newlabel{prop-FNLW@cref}{{[originalproposition][0][]5.1}{[1][0][]source}{}{}{}}
\makeatother
\crefname{originalproposition}{original proposition}{original propositions}
\Crefname{originalproposition}{Original proposition}{Original propositions}
\makeatletter
\newlabel{prop-glocal}{{4.1}{source}{}{sourcepointer-12}{}}
\newlabel{prop-glocal@cref}{{[originalproposition][0][]4.1}{[1][0][]source}{}{}{}}
\makeatother
\crefname{originaltheorem}{original theorem}{original theorems}
\Crefname{originaltheorem}{Original theorem}{Original theorems}
\makeatletter
\newlabel{thm-SDNLW}{{1.2}{source}{}{sourcepointer-13}{}}
\newlabel{thm-SDNLW@cref}{{[originaltheorem][0][]1.2}{[1][0][]source}{}{}{}}
\makeatother
\crefname{originaltheorem}{original theorem}{original theorems}
\Crefname{originaltheorem}{Original theorem}{Original theorems}
\makeatletter
\newlabel{thm-invariance}{{1.4}{source}{}{sourcepointer-14}{}}
\newlabel{thm-invariance@cref}{{[originaltheorem][0][]1.4}{[1][0][]source}{}{}{}}
\makeatother
\crefname{originalequation}{original equation}{original equations}
\Crefname{originalequation}{Original equation}{Original equations}
\makeatletter
\newlabel{trunc-SDNLW}{{1.14}{source}{}{sourcepointer-15}{}}
\newlabel{trunc-SDNLW@cref}{{[originalequation][0][]1.14}{[1][0][]source}{}{}{}}
\makeatother
\title{Author result prop-approx (printed 5.3) of ahl-163-stochastic-nonlinear-wave-dynamics-on-compact-2023}
\author{Oh, Tadahiro, Robert, Tristan, Tzvetkov, Nikolay}
\date{doi:10.5802/ahl.163; CC BY 4.0}
\begin{document}
\maketitle
\noindent\textit{Editorial note.} All delivered mathematical text is the author's own, copied from the pinned source sources/ahl-163-stochastic-nonlinear-wave-dynamics-on-compact-2023/pinned-original-source.tex.\par\medskip
\noindent\textit{Editorial note.} Local material id L42309, claim key doi:10.5802/ahl.163:prop-approx, lane-x3 S062 batch-2 assignment (CLAIMS six-line segment).\par\medskip
\noindent\textit{Editorial scope.} Single face of a Centre Mersenne journal unit: the counted text is the author's own statement and proof copied verbatim from the pinned source; the card adds no macro of its own beyond declaring the author environment names.\par\medskip
\noindent\textit{Allowed standard prerequisites (editorial).}\begin{itemize}
\item The statement and its own proof as printed in the article. source lines 2172-2369. the counted unit and the argument carried out in the span
\end{itemize}\medskip
\ifdefined\hypertarget\hypertarget{sourcepointer-1}{}\fi
\noindent\textit{Source reference.} Original equation 5.1: sources/ahl-163-stochastic-nonlinear-wave-dynamics-on-compact-2023/pinned-original-source.tex, source lines 1864--1864. Navigation pointer carrying the number the article prints (from census v3.14); nothing of the referenced item is copied.\par\medskip
\ifdefined\hypertarget\hypertarget{sourcepointer-2}{}\fi
\noindent\textit{Source reference.} Original equation 1.16: sources/ahl-163-stochastic-nonlinear-wave-dynamics-on-compact-2023/pinned-original-source.tex, source lines 380--380. Navigation pointer carrying the number the article prints (from census v3.14); nothing of the referenced item is copied.\par\medskip
\ifdefined\hypertarget\hypertarget{sourcepointer-3}{}\fi
\noindent\textit{Source reference.} Original equation 3.1: sources/ahl-163-stochastic-nonlinear-wave-dynamics-on-compact-2023/pinned-original-source.tex, source lines 894--894. Navigation pointer carrying the number the article prints (from census v3.14); nothing of the referenced item is copied.\par\medskip
\ifdefined\hypertarget\hypertarget{sourcepointer-4}{}\fi
\noindent\textit{Source reference.} Original equation 5.11: sources/ahl-163-stochastic-nonlinear-wave-dynamics-on-compact-2023/pinned-original-source.tex, source lines 2069--2069. Navigation pointer carrying the number the article prints (from census v3.14); nothing of the referenced item is copied.\par\medskip
\ifdefined\hypertarget\hypertarget{sourcepointer-5}{}\fi
\noindent\textit{Source reference.} Original equation 5.10: sources/ahl-163-stochastic-nonlinear-wave-dynamics-on-compact-2023/pinned-original-source.tex, source lines 2056--2056. Navigation pointer carrying the number the article prints (from census v3.14); nothing of the referenced item is copied.\par\medskip
\ifdefined\hypertarget\hypertarget{sourcepointer-6}{}\fi
\noindent\textit{Source reference.} Original equation 5.10: sources/ahl-163-stochastic-nonlinear-wave-dynamics-on-compact-2023/pinned-original-source.tex, source lines 2059--2059. Navigation pointer carrying the number the article prints (from census v3.14); nothing of the referenced item is copied.\par\medskip
\ifdefined\hypertarget\hypertarget{sourcepointer-7}{}\fi
\noindent\textit{Source reference.} Original equation 1.11: sources/ahl-163-stochastic-nonlinear-wave-dynamics-on-compact-2023/pinned-original-source.tex, source lines 318--318. Navigation pointer carrying the number the article prints (from census v3.14); nothing of the referenced item is copied.\par\medskip
\ifdefined\hypertarget\hypertarget{sourcepointer-8}{}\fi
\noindent\textit{Source reference.} Original equation 5.9: sources/ahl-163-stochastic-nonlinear-wave-dynamics-on-compact-2023/pinned-original-source.tex, source lines 2040--2040. Navigation pointer carrying the number the article prints (from census v3.14); nothing of the referenced item is copied.\par\medskip
\ifdefined\hypertarget\hypertarget{sourcepointer-9}{}\fi
\noindent\textit{Source reference.} Original equation 5.14: sources/ahl-163-stochastic-nonlinear-wave-dynamics-on-compact-2023/pinned-original-source.tex, source lines 2139--2139. Navigation pointer carrying the number the article prints (from census v3.14); nothing of the referenced item is copied.\par\medskip
\ifdefined\hypertarget\hypertarget{sourcepointer-10}{}\fi
\noindent\textit{Source reference.} Original lemma 3.2: sources/ahl-163-stochastic-nonlinear-wave-dynamics-on-compact-2023/pinned-original-source.tex, source lines 930--930. Navigation pointer carrying the number the article prints (from census v3.14); nothing of the referenced item is copied.\par\medskip
\ifdefined\hypertarget\hypertarget{sourcepointer-11}{}\fi
\noindent\textit{Source reference.} Original proposition 5.1: sources/ahl-163-stochastic-nonlinear-wave-dynamics-on-compact-2023/pinned-original-source.tex, source lines 1876--1876. Navigation pointer carrying the number the article prints (from census v3.14); nothing of the referenced item is copied.\par\medskip
\ifdefined\hypertarget\hypertarget{sourcepointer-12}{}\fi
\noindent\textit{Source reference.} Original proposition 4.1: sources/ahl-163-stochastic-nonlinear-wave-dynamics-on-compact-2023/pinned-original-source.tex, source lines 1625--1625. Navigation pointer carrying the number the article prints (from census v3.14); nothing of the referenced item is copied.\par\medskip
\ifdefined\hypertarget\hypertarget{sourcepointer-13}{}\fi
\noindent\textit{Source reference.} Original theorem 1.2: sources/ahl-163-stochastic-nonlinear-wave-dynamics-on-compact-2023/pinned-original-source.tex, source lines 350--350. Navigation pointer carrying the number the article prints (from census v3.14); nothing of the referenced item is copied.\par\medskip
\ifdefined\hypertarget\hypertarget{sourcepointer-14}{}\fi
\noindent\textit{Source reference.} Original theorem 1.4: sources/ahl-163-stochastic-nonlinear-wave-dynamics-on-compact-2023/pinned-original-source.tex, source lines 388--388. Navigation pointer carrying the number the article prints (from census v3.14); nothing of the referenced item is copied.\par\medskip
\ifdefined\hypertarget\hypertarget{sourcepointer-15}{}\fi
\noindent\textit{Source reference.} Original equation 1.14: sources/ahl-163-stochastic-nonlinear-wave-dynamics-on-compact-2023/pinned-original-source.tex, source lines 342--342. Navigation pointer carrying the number the article prints (from census v3.14); nothing of the referenced item is copied.\par\medskip
\begin{prop}\label{prop-approx}
Let $m,j\in\N$, $N_p\to\infty$ and $(u_0,u_1,\o) \in\bigcap_{p\in\N}\Si_{N_p}^{m,j}$. Then $\{\Phi^{N_p}(t)(u_0,u_1,\o)-\Psi(t)(u_0,u_1,\o)\}_{p\,\in\,\N}$ is a Cauchy sequence in $C([0,2^j];\H^{s_2}(\M))$. In particular, $\{\Phi^{N_p}(t)(u_0,u_1,\o)\}_{p\,\in\,\N}$ is a Cauchy sequence in $C([0,2^j];\H^s(\M))$. Here $s<0<s_2<s_1<1+s$ with $0<-s<1-s_1<1-s_2\ll 1$.
\end{prop}

\begin{proof}[Proof of Proposition~\ref{prop-approx}]
Let us fix $m,j\in\N$, $\dl>0$ as in~\eqref{dlmj}, and $(u_0,u_1,\o)\in\bigcap_{p\,\in\,\N}\Si_{N_p}^{m,j}$. In the following, we fix two (large) integers $N,M\in\{N_p\}_{p\in\N}$. Again, we write
\begin{equation*}
\left(w_N,\dt w_N\right)(t) = \Phi^N(t)\left(u_0,u_1,\o\right)-\Psi(t,u_0,u_1,\o),
\end{equation*}
and we denote by $\Phi^N_1(t)$ (respectively $\Phi^N_2(t)$) the first (respectively second) component of $\Phi^N(t)$. We will control inductively the difference $(w_N(t),\dt w_N(t))-(w_M(t),\dt w_M(t))$ on the time intervals $[a\dl,(a+1)\dl]$, $a=1,\,\dots,\,[2^j/\dl]$. We begin by controlling the difference on the first time interval, corresponding to $a=0$.

Then on $[0,\dl]$, we have for $N\le M$:
\begingroup
\allowdisplaybreaks
\begin{multline}\label{I-II-III}
w_{M}(t)-w_N(t)\\
\mkern-250mu\begin{aligned}
&= \sum_{\ell=0}^k{\binom k \ell}\int_0^tV(t-t')\bigg\{\P_{M}\Big[H_{\ell}\big(\P_{M}\Psi(t',u_0,u_1,\o);\s_{M}(x)\big)\big(\P_{M}w_{M}(t')\big)^{k-\ell}\Big]\\
&\qquad-\P_{N}\Big[H_{\ell}\big(\P_{N}\Psi(t',u_0,u_1,\o);\s_N(x)\big)\big(\P_Nw_N(t')\big)^{k-\ell}\Big]\bigg\}dt'
\\
&=\sum_{\ell=0}^{k}{\binom k \ell}\int_{0}^tV(t-t')\bigg\{(\P_{M}-\P_N)\Big[H_{\ell}\Big(\P_{M}\Psi(t',u_0,u_1,\o);\s_{M}(x)\Big)\big(\P_{M}w_{M}(t')\big)^{k-\ell}\Big]\\
&\qquad+\P_{N}\Big[\Big(H_{\ell}\big(\P_{M}\Psi(t',u_0,u_1,\o);\s_{M}(x)\big)\\
&\qquad-H_{\ell}\big(\P_{N}\Psi(t',u_0,u_1,\o);\s_N(x)\big)\Big)\big(\P_{M}w_{M}(t')\big)^{k-\ell}\Big]\\
&\qquad+\P_N\Big[H_{\ell}\big(\P_{N}\Psi(t',u_0,u_1,\o);\s_N(x)\big)\Big(\big(\P_{M}w_{M}(t')\big)^{k-\ell}-\big(\P_Nw_N(t')\big)^{k-\ell}\Big)\Big]\bigg\}dt'\\
&=:\I+\II+\III.
\end{aligned}
\end{multline}
\endgroup

To estimate these terms, we proceed as in the proof of Proposition~\ref{prop-glocal}. We begin with
\begingroup
\allowdisplaybreaks
\begin{multline*}
\|\I\|_{C\left([0,\dl];H^{s_2}\right)}\\
\begin{aligned}
&\les \sum_{\ell=0}^k\Big\|(\P_{M}-\P_N)\Big[H_{\ell}\Big(\P_{M}\Psi(u_0,u_1,\o);\s_{M}(x)\Big)\big(\P_{M}w_{M}\big)^{k-\ell}\Big]\Big\|_{L^1\left([0,\dl];H^{s_2-1}\right)}\\
&\les \sum_{\ell=0}^k N^{s_2-s_1}\Big\|\Big[H_{\ell}\Big(\P_{M}\Psi(u_0,u_1,\o);\s_{M}(x)\Big)\big(\P_{M}w_{M}\big)^{k-\ell}\Big]\Big\|_{L^1\left([0,\dl];H^{s_1-1}\right)}\\
&\les \sum_{\ell=0}^k N^{s_2-s_1}\dl^{\frac12}\left\|H_{\ell}\Big(\P_{M}\Psi(u_0,u_1,\o);\s_{M}(x)\Big)\right\|_{L^2\left([0,1];W^{s,\infty}\right)}\|\P_{M}w_{M}\|_{C\left([0,\dl];H^{s_1}\right)}^{k-\ell}\\
&\le C(m,j) N^{s_2-s_1},
\end{aligned}
\end{multline*}
\endgroup
for some constant $C(m,j)$ independent of $N,M$, where the second to last estimate comes from the same argument as in the proof of Proposition~\ref{prop-glocal}, and the last one from the condition~\eqref{condition2} given by $(u_0,u_1,\o)\in \B_{M}^{m,j}(D)$, with~\eqref{bd-wN} and the choice of $\dl$ in~\eqref{dlmj}.

Similarly, we bound
\begingroup
\allowdisplaybreaks
\begin{align*}
&\|\II\|_{C\left([0,\dl];H^{s_2}\right)}\\
&\les\sum_{\ell=1}^{k}\Big\|\Big[H_{\ell}\big(\P_{M}\Psi(u_0,u_1,\o);\s_{M}(x)\big)\\
&\qquad-H_{\ell}\big(\P_{N}\Psi(u_0,u_1,\o);\s_N(x)\big)\Big]\big(\P_{M}w_{M}\big)^{k-\ell}\Big\|_{L^1\left([0,\dl];H^{s_2-1}\right)}\\
&\les \dl^{\frac12}\bigg\{\Big\|H_{k}\big(\P_{M}\Psi(u_0,u_1,\o);\s_{M}(x)\big)-H_{k}\big(\P_{N}\Psi(u_0,u_1,\o);\s_N(x)\big)\Big\|_{L^2\left([0,\dl];H^{s_2-1}\right)}\\
&\qquad+\sum_{\ell=1}^{k-1}\Big\|H_{\ell}\big(\P_{M}\Psi(u_0,u_1,\o);\s_{M}(x)\big)-H_{\ell}\big(\P_{N}\Psi(u_0,u_1,\o);\s_N(x)\big)\Big\|_{L^2\left([0,\dl];W^{s,\infty}\right)}\\
&\qquad\qquad\times\big\|\P_{M}w_{M}\big\|_{C\left([0,\dl];H^{s_1}\right)}^{k-\ell}\bigg\}\\
& \le C(m,j) N^{-\eps},
\end{align*}
\endgroup
where the last bound comes from~\eqref{condition3} given by $(u_0,u_1,\o)\in\B_{M}^{m,j}(D)$ and from~\eqref{bd-wN} with the choice of $\dl$.

Finally, we can further decompose
%\begingroup
%\allowdisplaybreaks
\begin{multline*}
\III=\sum_{\ell=0}^{k-1}{\binom k \ell}\int_0^tV(t-t')\P_N\bigg\{H_{\ell}\big(\P_{N}\Psi(t',u_0,u_1,\o);\s_N(x)\big)\\
\begin{aligned}
&\Big[\Big(w_N(t')^{k-\ell}-\big(\P_Nw_N(t')\big)^{k-\ell}\Big)+\Big(\big(\P_{M}w_{M}(t')\big)^{k-\ell}-w_{M}(t')^{k-\ell}\Big)\\
&+\Big(w_{M}(t')^{k-\ell}-w_N(t')^{k-\ell}\Big)\Big]\bigg\}dt'=:\III_1+\III_2+\III_3.
\end{aligned}
\end{multline*}

We estimate similarly as before
\begin{align*}
\|\III_1\|_{C\left([0,\dl];H^{s_2}\right)}&\les \sum_{\ell=0}^{k-1}\dl^{\frac12}\big\|H_{\ell}\big(\P_{N}\Psi(u_0,u_1,\o);\s_N(x)\big)\big\|_{L^2([0,1];W^{s,\infty})}\\
&\quad\times\big\|(1-\P_{N})w_{N}\big\|_{C\left([0,\dl];H^{s_2}\right)}\left(\big\|\P_{N}w_{N}\big\|_{C\left([0,\dl];H^{s_2}\right)}^{k-\ell-1}+\big\|w_{N}\big\|_{C\left([0,\dl];H^{s_2}\right)}^{k-\ell-1}\right)\\
&\le C(m,j) N^{s_2-s_1},
\end{align*}
where the first estimate follows from the same argument as in the proof of Proposition~\ref{prop-glocal}, provided that $s_2<s_1<1+s<1$ is close enough to 1. The same argument applies to $\III_2$ and gives the same bound (with $M$ in place of $N$), and the last term can be bounded similarly by
\begingroup
\allowdisplaybreaks
\begin{align*}
\big\|\III_3\big\|_{C\left([0,\dl];H^{s_2}\right)}
&\les\sum_{\ell=0}^{k-1}\dl^{\frac12}\big\|H_{\ell}\big(\P_{N}\Psi(u_0,u_1,\o);\s_N(x)\big)\big\|_{L^2\left([0,1];W^{s,\infty}\right)}\\
&\quad\times\big\|w_{M}-w_N\big\|_{C\left([0,\dl];H^{s_2}\right)}\left(\big\|w_{M}\big\|_{C\left([0,\dl];H^{s_2}\right)}^{k-\ell-1}+\big\|w_{N}\big\|_{C\left([0,\dl];H^{s_2}\right)}^{k-\ell-1}\right)\\
&\le C\dl^{\frac12}D(m+j)^{\frac{k}2}\|w_{M}-w_N\|_{C([0,\dl];H^{s_2})},
\end{align*}
\endgroup
where the last estimate comes again from~\eqref{condition2} thanks to $(u_0,u_1,\o)\in\B_{M}^{m,j}(D)$, and from~\eqref{bd-wN} with the argument of Proposition~\ref{prop-glocal} applied with $s_2$ (provided that $s_2$ is sufficiently close to 1). With our choice of $C\dl^{\frac12}D(m+j)^{\frac{k}2}=C\dl^{\frac12}R\le \frac12$, we can absorb this last term in the left-hand side of~\eqref{I-II-III}.


The same arguments also apply to control $\Phi^N_2(t)-\Phi^M_2(t)$ on $[0,\dl]$. Therefore, gathering the estimates above leads to
\begin{equation}\label{S1}
\big\|\Phi^{M}(t)(u_0,u_1,\o)-\Phi^N(t)(u_0,u_1,\o)\big\|_{C\left([0,\dl];\H^{s_2}\right)} \le C(m,j)\big(N^{s_2-s_1}+N^{-\eps}\big)
\end{equation}
for any $N\le M$. In particular this shows the convergence on the time interval $[0,\dl]$.

We now investigate the convergence on the second time interval $[\dl,2\dl]$: we first decompose
\begin{multline*}
\left\|\Phi^{M}(t+\dl)(u_0,u_1,\o)-\Phi^N(t+\dl)(u_0,u_1,\o)\right\|_{C\left([0,\dl];\H^{s_2}\right)}\\
\begin{aligned}
& \le \left\|\Phi^{M}(t)\wt\Phi^{M}(\dl)(u_0,u_1,\o)-\Phi^N(t)\wt\Phi^{M}(\dl)(u_0,u_1,\o)\right\|_{C\left([0,\dl];\H^{s_2}\right)}\\
&\quad+\left\|\Phi^{N}(t)\wt\Phi^{M}(\dl)(u_0,u_1,\o)-\Phi^N(t)\wt\Phi^{N}(\dl)(u_0,u_1,\o)\right\|_{C\left([0,\dl];\H^{s_2}\right)}.
\end{aligned}
\end{multline*}

Note that replacing $\Psi(u_0,u_1,\o)$ by $\Psi(\wt\Phi^{M}(\dl)(u_0,u_1,\o))$ in the previous estimates and using that we still have $\wt\Phi^{M}(\dl)(u_0,u_1,\o)\in\B_{M}^{m,j}$ by choice of $(u_0,u_1,\o)\in\Si_{M}^{m,j}$ shows that the first term above is still bounded by
\begin{multline*}
\left\|\Phi^{M}(t)\wt\Phi^{M}(\dl)(u_0,u_1,\o)-\Phi^N(t)\wt\Phi^{M}(\dl)(u_0,u_1,\o)\right\|_{C\left([0,\dl];\H^{s_2}\right)}\\
\le C(m,j)\left(N^{s_2-s_1}+N^{-\eps}\right)
\end{multline*}
for any $N\le M$.

Thus we need to deal with the second term. We can redefine
\begin{align*}
w_{M}(t) &= \Phi_1^{N}(t)\wt\Phi^{M}(\dl)(u_0,u_1,\o) - \Psi\left(t,\wt\Phi^{M}(\dl)(u_0,u_1,\o)\right)
\\
\intertext{and}
w_{N}(t) &= \Phi_1^{N}(t)\wt\Phi^{N}(\dl)(u_0,u_1,\o) - \Psi\left(t,\wt\Phi^{N}(\dl)(u_0,u_1,\o)\right),
\end{align*}
and since $N,M\in\{N_p\}_{p\,\in\,\N}$ and $(u_0,u_1,\o)\in \bigcap_{p\,\in\,\N}\Si_{N_p}^{m,j}$ we have in particular that both $\wt\Phi^{M}(\dl)(u_0,u_1,\o)\in\B_M^{m,j}(D)$ and $\wt\Phi^{N}(\dl)(u_0,u_1,\o)\in\B_N^{m,j}(D)$ so that both $w_M$ and $w_N$ are well-defined and enjoy the bound~\eqref{bd-wN}. Moreover, by definition of $\Phi_1^N(t)$ and $\Psi$, they satisfy the following Duhamel formula:
\begin{align*}
w_M &= \sum_{\ell=0}^k{\binom k \ell}\int_0^tV(t-t')\P_N
\\
&\quad\bigg\{H_{\ell}\Big(\P_N\Psi\big(t',\wt\Phi^M(\dl)(u_0,u_1,\o)\big);\s_N(x)\Big)\big(\P_Nw_M(t')\big)^{k-\ell}\bigg\}dt',
\\
\intertext{and}
w_N &= \sum_{\ell=0}^k{\binom k \ell}\int_0^tV(t-t')\P_N\\
&\quad\bigg\{H_{\ell}\Big(\P_N\Psi\big(t',\wt\Phi^N(\dl)(u_0,u_1,\o)\big);\s_N(x)\Big)\big(\P_Nw_N(t')\big)^{k-\ell}\bigg\}dt'.
\end{align*}

To estimate in $C([0,\dl];H^{s_2}(\M))$ the difference
\begin{multline*}
\Phi_1^{N}(t)\wt\Phi^{M}(\dl)(u_0,u_1,\o)-\Phi_1^N(t)\wt\Phi^{N}(\dl)(u_0,u_1,\o)\\
=\Psi\big(t,\wt\Phi^{M}(\dl)(u_0,u_1,\o)\big)-\Psi\big(t,\wt\Phi^{N}(\dl)(u_0,u_1,\o)\big)+w_M-w_N,
\end{multline*}
we first bound directly the linear terms by
\begin{multline*}
\left\|\Psi\left(t,\wt\Phi^{M}(\dl)(u_0,u_1,\o)\right)-\Psi\left(t,\wt\Phi^{N}(\dl)(u_0,u_1,\o)\right)\right\|_{C\left([0,\dl];H^{s_2}\right)}\\
\begin{aligned}
&\le \left\|\Phi^{M}(\dl)(u_0,u_1,\o)-\Phi^{N}(\dl)(u_0,u_1,\o)\right\|_{\H^{s_2}}\\
& \le C(m,j)\left(N^{s_2-s_1}+N^{-\eps}\right)
\end{aligned}
\end{multline*}
thanks to~\eqref{def-Psidamp} and~\eqref{S1}.

To estimate the difference of the nonlinear components, we decompose
\begin{align*}
&w_M-w_N\\
&= \sum_{\ell=0}^k{\binom k \ell}\int_0^tV(t-t')\P_N\bigg\{H_{\ell}\Big(\P_N\Psi\left(t',\wt\Phi^M(\dl)(u_0,u_1,\o)\right);\s_N(x)\Big)\\
&\quad\big(\P_Nw_M(t')\big)^{k-\ell}
-H_{\ell}\Big(\P_N\Psi\left(t',\wt\Phi^N(\dl)(u_0,u_1,\o)\right);\s_N(x)\Big)\left(\P_Nw_N(t')\right)^{k-\ell}\bigg\}dt'\\
&=\sum_{\ell=0}^k{\binom k \ell}\int_0^tV(t-t')\P_N\bigg\{\Big[H_{\ell}\Big(\P_N\Psi\big(t',\wt\Phi^M(\dl)(u_0,u_1,\o)\big);\s_N(x)\Big)\\
&\quad-H_{\ell}\Big(\P_N\Psi\left(t',\wt\Phi^N(\dl)(u_0,u_1,\o)\right);\s_N(x)\Big)\Big]\big(\P_Nw_M(t')\big)^{k-\ell}\\
&\quad+H_{\ell}\Big(\P_N\Psi\left(t',\wt\Phi^N(\dl)(u_0,u_1,\o)\right);\s_N(x)\Big)\\
&\mkern200mu\Big[\big(\P_Nw_M(t')\big)^{k-\ell}
-\big(\P_Nw_N(t')\big)^{k-\ell}\Big]\bigg\}dt'\\
&=:\wt \II + \wt\III.
\end{align*}

First note that we can estimate $\wt \III$ exactly as $\III$ in~\eqref{I-II-III}, giving the bound
\begin{equation*}
\big\|\wt\III\big\|_{C\left([0,\dl];H^{s_2}\right)}\le C(m,j)N^{s_2-s_1} + \frac12\big\|w_M-w_N\big\|_{C\left([0,\dl];H^{s_2}\right)}.
\end{equation*}

Finally, we estimate the remaining term by
\begin{multline*}
\big\|\wt\II\big\|_{C\left([0,\dl];H^{s_2}\right)}\les \sum_{\ell=1}^k\Big\|\Big[H_{\ell}\Big(\P_N\Psi\left(\wt\Phi^M(\dl)(u_0,u_1,\o)\right);\s_N(x)\Big)\\
-H_{\ell}\Big(\P_N\Psi\left(\wt\Phi^N(\dl)(u_0,u_1,\o)\right);\s_N(x)\Big)\Big]\big(\P_Nw_M\big)^{k-\ell}\Big\|_{L^1\left([0,\dl];H^{s_2-1}\right)}
\end{multline*}
Writing then
\begin{align*}
H_{\ell}(u;\s_N)-H_{\ell}(v;\s_N) = -\sum_{i=0}^{\ell-1}{\binom \ell i}H_i(u;\s_N)(v-u)^{\ell-i}
\end{align*}
thanks to~\eqref{Hsum}, we can then estimate the previous term with
\begin{multline*}
\sum_{\ell=1}^k\sum_{i=0}^{\ell-1}\dl^{\frac12}\Big\|\Big[\P_N\Psi\left(\wt\Phi^M(\dl)(u_0,u_1,\o)\right)-\P_N\Psi\left(\wt\Phi^N(\dl)(u_0,u_1,\o)\right)\Big]^{\ell-i}\\
\begin{aligned}
&\quad\times H_{i}\Big(\P_N\Psi\left(\wt\Phi^N(\dl)(u_0,u_1,\o)\right);\s_N(x)\Big)\big(\P_Nw_M\big)^{k-\ell}\Big\|_{L^2\left([0,\dl];H^{s_2-1}\right)}\\
&\les \sum_{\ell=1}^k\sum_{i=0}^{\ell-1}\dl^{\frac12}\Big\|H_i\Big(\P_N\Psi\left(\wt\Phi^N(\dl)(u_0,u_1,\o)\right);\s_N(x)\Big)\Big\|_{L^2\left([0,1];W^{s,\infty}\right)}\\
&\quad\times\Big\|\Psi\left(\wt\Phi^M(\dl)(u_0,u_1,\o)\right)-\Psi\left(\wt\Phi^N(\dl)(u_0,u_1,\o)\right)\Big\|_{C\left([0,\dl];H^{s_2}\right)}^{\ell-i}\big\|w_M\big\|_{C\left([0,\dl];H^{s_2}\right)}^{k-\ell},
\end{aligned}
\end{multline*}
provided again that $s_2$ is close enough to 1 (depending on $k$). Using then that $\wt\Phi^N(\dl)(u_0,u_1,\o)\in\B_N^{m,j}(D)$ and~\eqref{bd-wN},\eqref{S1}, we finally get
\begin{equation*}
\big\|\wt\II\big\|_{C\left([0,\dl];H^{s_2}\right)}\le C(m,j)\left(N^{s_2-s_1}+N^{-\eps}\right).
\end{equation*}

Gathering the estimates above, we obtain
\begin{align*}
\big\|w_M-w_N\big\|_{C\left([0,\dl];H^{s_2}\right)}\le 4C(m,j)\left(N^{s_2-s_1}+N^{-\eps}\right)
\end{align*}
which leads to
\begin{align*}
\big\|(w_N,\dt w_N)(t+\dl)-(w_M,\dt w_M)(t+\dl)\big\|_{C\left([0,\dl];\H^{s_2}\right)} \le C_2(m,j)\left(N^{s_2-s_1}+N^{-\eps}\right)
\end{align*}
for some larger constant $C_2(m,j)\ge C(m,j)$. This shows that $\{(w_{N_p},\dt w_{N_p})\}_{p\,\in\,\N}$ is also a Cauchy sequence in $C([0,2\dl];\H^{s_2}(\M))$.

We can then proceed inductively on $a=0,\,\dots,\,[2^j/\dl]$ and repeat the previous estimates by using that at each step $\wt\Phi^N(a\dl)(u_0,u_1,\o)\in\B_N^{m,j}(D)$ since $(u_0,u_1,\o)\in \Si_N^{m,j}$. Thus we deduce that there exists a (large) constant $C_{2^j/\dl}(m,j)>0$ such that for any $N,M \in \{N_p\}_{p\,\in\,\N}$ with $N\le M$ it holds
\begin{equation}
\big\|(w_N, \dt w_N)-(w_M,\dt w_M)\big\|_{C\left([0,2^j];\H^{s_2}\right)} \le C_{2^j/\dl}(m,j)\left(N^{s_2-s_1}+N^{-\eps}\right).
\end{equation}
This is enough to show the convergence of $\{\Phi^{N_p}(t)(u_0,u_1,\o)-\Psi(t)(u_0,u_1,\o)\}$ in $C([0,2^j];\H^{s_2}(\M))$. As a result, $\big\{\Phi^{N_p}(t)(u_0,u_1,\o)\big\}_{p\,\in\,\N}$ converges in $C([0,2^j];\linebreak\H^s(\M))$. This concludes the proof of Proposition~\ref{prop-approx}.
\end{proof}

\end{document}
